\section{Referencial teórico}\label{cap:refteo}

Este capítulo constrói a lente analítica utilizada para leitura do corpus. Partimos do debate de longa duração sobre máquinas e emprego, passamos pelo paradigma da mudança técnica enviesada por qualificação e por seu limite empírico -- a polarização do trabalho --, chegamos ao arcabouço de tarefas de \textcite{acemogluRestrepo2019}, cujos quatro mecanismos estruturam a codificação dos estudos, e termina situando a literatura específica da inteligência artificial e a virada generativa. O objetivo não é esgotar a teoria econômica do progresso técnico, mas fixar com precisão os conceitos -- deslocamento, reinstalação, complementaridade e demanda agregada -- que reaparecem, já operacionalizados, nos Capítulos \ref{cap:janelas} e \ref{cap:comparacao}, e explicitar a lacuna que a revisão se propõe a examinar.

\subsection{Tecnologia e emprego: um debate de longa duração}

A preocupação de que as máquinas suprimam o trabalho humano é mais antiga do que a economia como disciplina formal, mas foi \textcite{keynes1930} quem popularizou o conceito que ainda hoje organiza o debate: o \emph{desemprego tecnológico}, isto é, o desemprego decorrente de que a economia descobre como economizar mão de obra num ritmo mais rápido do que o da abertura de novos usos para o trabalho.  Keynes tratava o fenômeno como uma disfunção transitória de um processo, no  fundo, benigno, qual seja, a elevação secular do padrão de vida. Essa tensão entre  um efeito-deslocamento de curto prazo e um efeito-compensação de prazo mais longo atravessa toda a literatura posterior.

A versão mais célebre do alarme veio de \textcite{leontief1983}, para quem o papel do trabalho humano tenderia a encolher como encolhera o papel do cavalo na agricultura. A analogia representa a hipótese de que, a partir de certo ponto, a substituição deixaria de ser compensada. A história econômica, contudo, não confirmou a previsão pessimista. Como sintetiza \textcite{autor2015}, ao longo do século XX a automação eliminou tarefas em larga escala sem produzir desemprego tecnológico persistente, porque operou em paralelo a forças que
ampliaram a demanda por trabalho, isto é, a complementaridade entre máquinas e trabalhadores nas tarefas não automatizadas, o aumento de produtividade e renda e a criação de novas ocupações. O problema teórico relevante, portanto, não foi se a tecnologia destrói postos, mas qual o saldo entre destruição e criação e como ele se distribui entre os trabalhadores. É essa indeterminação que o restante do capítulo procura organizar conceitualmente, e que os resultados encontram, no agregado, como um sinal predominantemente ambíguo sobre o emprego (Capítulo~\ref{cap:janelas}).

% \subsection{Mudança técnica enviesada por qualificação e seu limite empírico}

A primeira resposta sistemática da economia do trabalho a essa indeterminação foi a hipótese da \emph{mudança técnica enviesada por qualificação}, ou \gls{sbtc}. Sua intuição remonta à ideia, formalizada por \textcite{tinbergen1974, tinbergen1975}, de uma ``corrida entre educação e tecnologia'': o progresso técnico elevaria continuamente a demanda relativa por trabalhadores qualificados, e o prêmio salarial da qualificação dependeria de a oferta de trabalhadores instruídos acompanhar ou não esse ritmo. \textcite{katzMurphy1992} deram à tese sua forma empírica canônica, explicando a elevação do prêmio universitário nos Estados Unidos a partir dos anos 1980 pela interação entre uma demanda relativa em alta -- atribuída à tecnologia, em especial à informática, hipótese que fora levantada por \textcite{krueger1991} -- e oscilações na oferta de formados. Os autores observaram que o prêmio salarial da qualificação subia mesmo com o aumento do número de graduados no mercado. \textcite{goldinKatz2008} estenderam a análise da corrida entre educação e tecnologia para todo o século XX da economia americana.

\textcite{acemoglu2002} consolidou o arcabouço teórico do SBTC, formalizando a ideia de que o progresso técnico pode ser enviesado por qualificação, ou seja, favorecer mais a mão de obra qualificada do que a não qualificada. O modelo de Acemoglu é o ponto de partida para a análise da mudança técnica enviesada por qualificação e seus efeitos sobre a distribuição de renda e emprego. O SBTC se tornou o paradigma dominante para explicar o aumento da desigualdade salarial e a elevação do prêmio da qualificação nas últimas décadas do século XX.

O poder explicativo do SBTC, porém, esbarrou num fato que ele não acomodava bem. A partir dos anos 1990, a estrutura ocupacional de diversas economias desenvolvidas não se deslocou de forma monotônica em favor da alta qualificação, e sim se \emph{polarizou}: cresceram simultaneamente as ocupações de alta remuneração e as de baixa remuneração, enquanto encolhiam as ocupações de remuneração intermediária. \textcite{autorKatzKearney2006, autorKatzKearney2008} documentaram esse padrão para os Estados Unidos e \textcite{goosManning2007} para o Reino Unido -- os ``\textit{lousy and lovely jobs}''. Um modelo em que a tecnologia simplesmente favorece ``mais qualificação'' não prevê o esvaziamento do \emph{meio} da distribuição; era preciso uma teoria que distinguisse não trabalhadores por nível de instrução, mas \emph{tarefas} por seu conteúdo.

% \subsection{A abordagem de tarefas e a polarização do trabalho}

Essa teoria veio com a \emph{abordagem de tarefas}. \textcite{autorLevyMurnane2003} propuseram deslocar a unidade de análise do trabalhador para a tarefa e classificaram as tarefas segundo serem ou não \emph{rotineiras}, isto é passíveis de descrição por regras explícitas. Os autores, no que ficou conhecido, também, como ``modelo ALM'' -- em referência às iniciais dos autores --, propuseram a distinção conceitual entre ``habilidades'' e ``tarefas''. As habilidades são atributos do trabalhador, como escolaridade ou experiência, e as tarefas são as unidades de atividade produtiva que precisam ser executadas para gerar um bem ou serviço, como digitar ou interpretar um texto. Os autores defendem que a tecnologia e a computação, em particular, pouco dizem respeito ao estoque de capacidades humanas do trabalhador -- suas habilidades --, mas sim ao conteúdo da rotina das atividades. Dessa forma, o modelo ALM propõe uma matriz de classificação das tarefas:

\begin{enumerate}
  \item Rotineiras Cognitivas: tarefas de escritório que seguem regras explícitas de processamento de dados (por exemplo, escriturários e contabilidade básica);
  \item Rotineiras Manuais: atividades fabris repetitivas ou linhas de produção (por exemplo, operários de montagem);
  \item Não Rotineiras Cognitivas (ou Abastratas): resolução de problemas complexos, gestão, engenharia, criação (por exemplo, gerentes, cientistas);
  \item Não Rotineiras Manuais: atividades que exigem habilidades físicas e manipulação de equipamentos, ou seja, adaptação motora ao ambiente (por exemplo, cuidadores de idosos, garçons).
\end{enumerate}

O modelo ALM formulou a hipótese de que a automação, ao substituir tarefas rotineiras, reduziria a demanda por trabalhadores de qualificação intermediária, enquanto os trabalhadores de alta qualificação, com tarefas não rotineiras analíticas e interativas, e os de baixa qualificação, com tarefas manuais não rotineiras, seriam relativamente poupados. O modelo explicava a polarização do trabalho como resultado da substituição de tarefas rotineiras por tecnologia, sem necessariamente favorecer mais os trabalhadores de alta qualificação do que os de baixa qualificação.


A tecnologia computacional, nessa leitura, substitui de forma barata as tarefas rotineiras -- tanto manuais quanto cognitivas -- e \emph{complementa} as tarefas não rotineiras, sejam as analíticas e interativas de alta qualificação, sejam as manuais não rotineiras de baixa qualificação, difíceis de codificar. Como as tarefas rotineiras se concentravam justamente nas ocupações de remuneração intermediária, a hipótese da \emph{mudança técnica enviesada por rotina}, ou \gls{rbtc}, explicava de uma só vez por que o meio se esvazia e as duas pontas crescem.

A evidência empírica subsequente sustentou amplamente o mecanismo. \textcite{autorDorn2013} mostraram que as áreas dos Estados Unidos historicamente mais intensivas em tarefas rotineiras foram as que mais se polarizaram, com expansão de empregos de serviços de baixa qualificação, e \textcite{goosManningSalomons2014} estenderam o resultado a dezesseis países europeus, atribuindo a polarização à combinação de RBTC e \textit{offshoring}. A síntese de referência do programa é \textcite{acemogluAutor2011}, que formaliza o ``\textit{task framework}'' e articula qualificações, tarefas e tecnologias num mesmo modelo de oferta e demanda. É dessa tradição que provém o conceito de \emph{polarização} tal como a revisão o codifica: a tese de que o risco da automação recai de forma desigual, concentrando-se nas tarefas rotineiras de menor qualificação. Como mostram os Capítulos \ref{cap:janelas} e \ref{cap:comparacao}, essa continua a ser a categoria modal de polarização no corpus, antes e depois da IA generativa.

\subsection{O arcabouço Acemoglu--Restrepo: deslocamento, reinstalação e a corrida entre tarefas}

A abordagem de tarefas ganhou sua formulação dinâmica mais influente na série de trabalhos de Daron Acemoglu e Pascual Restrepo, que constitui a espinha dorsal analítica desta revisão. O ponto de partida é distinguir, dentro do progresso técnico, forças com sinais opostos sobre a demanda por trabalho \parencite{acemogluRestrepo2018, acemogluRestrepo2019}. São quatro os mecanismos que a revisão extrai desse arcabouço e codifica em cada estudo:

\begin{description}
  \item[Deslocamento (\textit{displacement}):] O efeito direto da automação. Quando o capital passa a executar tarefas antes realizadas por trabalhadores, reduz-se a demanda por trabalho e a participação do trabalho na renda. É o canal que captura a "substituição" no sentido clássico e, como mostram os resultados, o mecanismo mais invocado pela literatura em todas as janelas.
  \item[Reinstalação (\textit{reinstatement}):] A força contrária de mesma natureza. A criação de \emph{novas} tarefas em que o trabalho tem vantagem comparativa. Para \textcite{acemogluRestrepo2019}, é sobretudo o surgimento de novas tarefas -- e não apenas o barateamento das antigas -- que sustenta a demanda por trabalho no longo prazo. A "corrida" do título de \textcite{acemogluRestrepo2018} é precisamente entre deslocamento e reinstalação.
  \item[Complementaridade e produtividade:] Mesmo nas tarefas que automatiza, a tecnologia pode baratear a produção, elevar a produtividade e, por essa via, expandir a demanda por trabalho nas tarefas não automatizadas. É o canal pelo qual a IA "amplia capacidades" em vez de apenas substituir, e o que mais cresce na literatura recente (Capítulo~\ref{cap:comparacao}).
  \item[Demanda agregada:] O efeito de renda em nível macroeconômico. Ganhos de produtividade que se convertem em maior renda e gasto e, consequentemente, em demanda por trabalho em toda a economia, inclusive em setores não diretamente afetados pela tecnologia. \textcite{bessen2018} ressalta que a elasticidade da demanda do setor é decisiva nesse canal: onde a demanda é elástica, a automação que barateia a produção pode \emph{elevar} o emprego setorial, e apenas onde a demanda já está saturada o efeito-deslocamento prevalece.
\end{description}

A virtude do arcabouço é tornar o saldo sobre o emprego uma questão \emph{empírica} sobre a magnitude relativa dessas forças, e não um veredito teórico fixo. Quando o deslocamento supera a reinstalação e os efeitos de produtividade, a participação do trabalho cai e o emprego se contrai; no caso oposto, a tecnologia é trabalho-amistosa. \textcite{acemogluRestrepo2020robots} estimam, para a robótica industrial nos Estados Unidos, um efeito líquido negativo sobre emprego e salários locais, ilustrando que o deslocamento pode, de fato, predominar; \textcite{acemogluRestrepo2022} mostram que boa parte do aumento da desigualdade salarial americana nas últimas décadas é compatível com um deslocamento de tarefas pouco compensado por reinstalação. Acemoglu e Restrepo advertem, por fim, que a \emph{direção} do progresso técnico é endógena: uma automação que apenas corta custos de mão de obra sem criar tarefas -- o que chamam de "\textit{so-so automation}" -- gera o pior dos mundos para o trabalho, deslocando sem reinstalar nem elevar muito a produtividade \parencite{acemogluRestrepo2020wrongai}. Essa moldura -- quatro forças cujo balanço determina o resultado -- é a que a revisão aplica, estudo a estudo, para
classificar como a literatura concebe o efeito da IA sobre o trabalho.

% \subsection{Antecedentes pré-2013 e a fronteira da revisão}

A revisão fixa o início de sua janela de busca em 2013, e a escolha tem justificativa teórica e empírica. Às vésperas desse marco, a economia do trabalho havia consolidado o arcabouço de tarefas como linguagem dominante e acumulado evidência robusta sobre a automação \emph{clássica} -- a informática e a robótica industrial. \textcite{graetzMichaels2018}, analisando dados de dezessete países, estimaram que os robôs industriais elevaram a produtividade sem reduzir o emprego agregado, ao passo que \textcite{acemogluRestrepo2020robots} chegaram a um saldo local negativo nos Estados Unidos; a coexistência dessas duas leituras define bem o estado da arte herdado. No plano da síntese de amplo alcance, \textcite{brynjolfssonMcafee2014} popularizaram a tese de uma "segunda era das máquinas" em que a digitalização passaria a afetar também o trabalho cognitivo, e \textcite{autor2015} ofereceu o contraponto cauteloso, lembrando a persistência histórica do emprego e o papel da complementaridade.

O ano de 2013 marca, ainda, a fronteira a partir da qual o objeto da revisão deixa de ser a automação genérica e passa a incluir a \emph{inteligência artificial} como categoria própria de tecnologia. Não por acaso, é de 2013 a primeira circulação do estudo de Frey e Osborne sobre a suscetibilidade das ocupações à computadorização, discutido adiante. A janela 2013--2025, portanto, não é arbitrária: ela cobre o período em que a IA -- primeiro o aprendizado de máquina e o \textit{deep learning}, depois os modelos generativos -- se torna o foco explícito da literatura, sem perder de vista o estoque conceitual herdado da era da automação. A justificativa operacional desse recorte é detalhada na metodologia (Capítulo~\ref{cap:metodologia}); aqui importa registrar que a fronteira temporal acompanha uma fronteira tecnológica.

% \subsection{Da automação à inteligência artificial: exposição ocupacional e a virada generativa}

A transposição do arcabouço de tarefas para a inteligência artificial enfrenta uma dificuldade específica: ao contrário da automação clássica, cujas tarefas-alvo eram visivelmente rotineiras, a IA contemporânea avança sobre tarefas \emph{não rotineiras} e cognitivas que se supunham protegidas. O esforço de medir essa exposição organiza boa parte da literatura empírica do período. O marco inaugural é \textcite{frey2017}, que estimou serem cerca de 47\% dos empregos nos Estados Unidos de alto risco de computadorização nas duas décadas seguintes -- número amplamente debatido e criticado por tratar ocupações inteiras, e não tarefas, como unidade de risco. As medidas posteriores refinaram justamente esse ponto. \textcite{brynjolfssonMitchellRock2018} propuseram avaliar a "\textit{suitability for machine learning}" tarefa a tarefa; \textcite{felten2018, felten2021} construíram o índice de exposição ocupacional à IA que os capítulos descritivos identificam no corpus, ligando avanços em capacidades de IA a habilidades ocupacionais; e \textcite{webb2020} usou o texto de patentes para mapear quais tarefas a IA, ao contrário de softwares e robôs anteriores, tende a atingir.

Um achado recorrente dessa geração de estudos é que o perfil de exposição da IA difere do da automação clássica: enquanto robôs e softwares atingiam tarefas rotineiras de qualificação intermediária e baixa, a exposição à IA preditiva tende a alcançar tarefas de qualificação mais alta, invertendo parcialmente o gradiente da RBTC \parencite{webb2020, felten2021}. Em evidência de demanda por trabalho, \textcite{acemogluVacancies2022} examinaram vagas \textit{online} e encontraram, nos estabelecimentos mais expostos à IA, recomposição de tarefas e reformulação de requisitos, sem efeito agregado nítido sobre o emprego até então -- um padrão coerente com a ambivalência que a revisão documenta.

A virada decisiva, contudo, é a da \emph{IA generativa}, cuja difusão a partir do lançamento público do ChatGPT em novembro de 2022 organiza a comparação central desta revisão (Capítulo~\ref{cap:comparacao}). Diferentemente das ondas anteriores, a evidência inicial sobre os modelos generativos é fortemente \emph{experimental} e centrada em produtividade. \textcite{noyZhang2023}, em experimento controlado de tarefas de escrita profissional, e \textcite{brynjolfssonLiRaymond2025}, com dados de atendimento ao cliente em firma real, encontraram ganhos de produtividade concentrados nos trabalhadores \emph{menos} experientes ou de menor desempenho inicial -- um efeito de \emph{compressão} que sugere complementaridade, e não apenas substituição. \textcite{eloundou2024}, por sua vez, estimaram que os grandes modelos de linguagem têm potencial de afetar parcela substancial das tarefas de \emph{muitas} ocupações, inclusive as de alta remuneração e escolaridade, tratando os \glspl{llm} como uma tecnologia de propósito geral. Trabalhos como o de \textcite{korinek2023} ilustram, do lado da própria pesquisa econômica, essa penetração da IA em tarefas cognitivas qualificadas. É essa migração do risco -- da tarefa rotineira de baixa qualificação para a tarefa cognitiva de alta qualificação -- que constitui a hipótese substantiva examinada pela revisão.

A interpretação teórica dessa evidência inicial está longe de ser consensual, e três leituras disputam o seu sentido. A primeira é cautelosa quanto ao agregado: \textcite{acemoglu2025simple} estima que, dadas a fração de tarefas efetivamente expostas e as economias de custo observadas, os ganhos macroeconômicos da IA na próxima década serão modestos, e \textcite{brynjolfsson2022turing} adverte que uma IA concebida para imitar o trabalho humano, em vez de ampliá-lo, tende a substituir sem criar -- a "armadilha de Turing", algo similiar à ideia de \textit{so-so automation}. A segunda leitura desloca a pergunta da substituição para a natureza do que é automatizado: \textcite{agrawalGansGoldfarb2019} tratam a IA como um barateamento da \emph{predição}, cujo efeito sobre o trabalho é ambíguo porque a predição tanto substitui tarefas quanto eleva o valor do julgamento humano que a complementa. A terceira é a mais otimista e a que mais diretamente dialoga com a evidência de compressão: para \textcite{autor2024rebuild}, ao baratear a \emph{expertise}, a IA generativa poderia permitir que trabalhadores de qualificação intermediária executem tarefas antes reservadas a especialistas, revertendo parte da polarização. Essas leituras importam para a revisão por duas razões. Tornam visível que complementaridade e reinstalação não são sinônimos -- ganhar produtividade nas tarefas existentes não é o mesmo que criar tarefas novas, e só a segunda sustenta a demanda por trabalho no longo prazo --, distinção que os resultados mostrarão ser empiricamente relevante. E sugerem que a dicotomia entre alta e baixa qualificação, herdada do debate SBTC-RBTC, pode ser estreita para uma tecnologia cujos efeitos mais bem documentados ocorrem \emph{dentro} das ocupações.

% \subsection{Síntese e lacuna: a hipótese da revisão}

O percurso teórico deste capítulo pode ser resumido em três movimentos. O SBTC explicou o prêmio da qualificação mas não a polarização; a abordagem de tarefas explicou a polarização ao deslocar o foco do trabalhador para a tarefa rotineira; e o arcabouço Acemoglu-Restrepo tornou o saldo sobre o emprego uma questão do balanço entre deslocamento, reinstalação, complementaridade e demanda agregada. Sobre esse acúmulo, a literatura específica da IA acrescentou uma conjectura nova e empiricamente aberta: a de que a inteligência artificial -- e, sobretudo, a generativa - romperia o padrão herdado da automação ao deslocar o risco para as tarefas cognitivas de \emph{alta} qualificação, antes tidas como complementares à máquina.

Daí decorre a lacuna que a revisão se propõe a examinar. A literatura econômica sobre IA e trabalho cresceu de forma explosiva e dispersa, com desenhos heterogêneos e vereditos não convergentes; falta uma síntese sistemática que verifique se o diagnóstico mudou \emph{de fato} com a chegada da IA generativa ou se a moldura da polarização e do efeito-substituição, herdada da era da automação, sobreviveu à transição. Em termos operacionais, a hipótese de trabalho sustenta que a literatura pré-ChatGPT convergia para a polarização com risco concentrado nas tarefas rotineiras de baixa qualificação, ao passo que a literatura pós-ChatGPT deslocaria esse risco para as tarefas cognitivas de alta qualificação -- uma hipótese de \emph{ruptura} testada, dimensão a dimensão, no capítulos de resultados. Os quatro mecanismos definidos n Seção~\ref{cap:refteo} fornecem as categorias dessa codificação; a comparação pareada pré/pós-ChatGPT (Capítulo~\ref{cap:comparacao}) fornece o teste. A própria abordagem de tarefas oferece, contudo, uma previsão mais fina do que a simples oposição entre períodos: se o risco segue o conteúdo das tarefas que cada tecnologia executa, o deslocamento do diagnóstico deve aparecer nos estudos cujo \emph{objeto} é a IA generativa, e não nos que seguem estudando robôs e automação, qualquer que seja a data em que foram publicados. Essa implicação é testada no Capítulo~\ref{cap:discussao}, usando a literatura de robótica como grupo de comparação. O que se verificará é menos uma inversão do consenso anterior do que um deslocamento parcial do diagnóstico, nítido onde a teoria o prevê e ausente onde ela não o prevê.
